% I may change the way this is done in a future version, 
%  but given that some people needed it, if you need a different degree title 
%  (e.g. Master of Science, Master in Science, Master of Arts, etc)
%  uncomment the following 3 lines and set as appropriate (this *has* to be before \maketitle)
% \makeatletter
% \renewcommand {\@degree@string} {Master of Things}
% \makeatother

\title{Exploiting Structural Synergy in post-SIRT imaging: Anatomically Guided Joint Variational Reconstruction of Yttrium-90 PET/SPECT/CT}
\author{Sam Porter}
\department{The institute of Nuclear Medicine, Division of Medicine}

\maketitle
\makedeclaration

\begin{abstract} % 300 word limit
Selective Internal Radiotherapy (SIRT) uses Yttrium-90 ($^{90}Y$) microspheres to treat unresectable liver tumours. Accurate post-treatment imaging is required to verify microsphere distribution and perform dosimetric evaluations. However, imaging with $^{90}Y$ presents unique challenges. Bremsstrahlung SPECT offers high count statistics but suffers from poor spatial resolution and difficulties in correcting for scatter and other additive terms. $^{90}Y$ PET offers superior resolution, but suffers from high noise due to the extremely low internal pair branching ratio ($\sim 0.003\%$). Traditional reconstruction fails to exploit complementary information in these paired datasets.

This thesis proposes a synergistic framework that uses anatomically guided joint variational regularisation to exploit shared underlying structure. We introduce the Directional Total Nuclear Variation Prior, which promotes edge alignment between PET and SPECT without enforcing intensity correlation, while incorporating CT edge-guidance and suppressing unwanted noise.

Further to this, we introduce a simple method of bremsstrahlung range modelling approach and residual correction routine to improve quantification. 

To ensure computational efficiency, we propose a synergy-aware preconditioner which, in conjunction with a variance reduced stochastic gradient algorithms, results in fast convergence.

We validate our method against the state-of-the-art sequential hybrid kernel expectation maximisation scheme (HKEM). We demonstrate that dTNV shows improved quantitative accuracy and increased reproducibility over noise realisations compared to existing methods.
\end{abstract}

\begin{impactstatement}

	UCL theses now have to include an impact statement. \textit{(I think for REF reasons?)} The following text is the description from the guide linked from the formatting and submission website of what that involves. (Link to the guide: {\scriptsize \url{http://www.grad.ucl.ac.uk/essinfo/docs/Impact-Statement-Guidance-Notes-for-Research-Students-and-Supervisors.pdf}})

\begin{quote}
The statement should describe, in no more than 500 words, how the expertise, knowledge, analysis,
discovery or insight presented in your thesis could be put to a beneficial use. Consider benefits both
inside and outside academia and the ways in which these benefits could be brought about.

The benefits inside academia could be to the discipline and future scholarship, research methods or
methodology, the curriculum; they might be within your research area and potentially within other
research areas.

The benefits outside academia could occur to commercial activity, social enterprise, professional
practice, clinical use, public health, public policy design, public service delivery, laws, public
discourse, culture, the quality of the environment or quality of life.

The impact could occur locally, regionally, nationally or internationally, to individuals, communities or
organisations and could be immediate or occur incrementally, in the context of a broader field of
research, over many years, decades or longer.

Impact could be brought about through disseminating outputs (either in scholarly journals or
elsewhere such as specialist or mainstream media), education, public engagement, translational
research, commercial and social enterprise activity, engaging with public policy makers and public
service delivery practitioners, influencing ministers, collaborating with academics and non-academics
etc.

Further information including a searchable list of hundreds of examples of UCL impact outside of
academia please see \url{https://www.ucl.ac.uk/impact/}. For thousands more examples, please see
\url{http://results.ref.ac.uk/Results/SelectUoa}.
\end{quote}
\end{impactstatement}

\begin{acknowledgements}
So, it turns out that a PhD is quite an undertaking. It's been a long old journey, but the end is in sight!

I've had an extraordinary four years and would like to thank some extraordinary people. Firstly, I'd like to thank my supervisors, Kris, Daniel, and Simon, who have been a near-bottomless source of knowledge, inspiration and ideas. Secondly, my colleagues in the INM, especially my fellow students. This would have been a very difficult (and much lonelier) few years without you all. I'm so grateful for our conversations in and out of the workplace. I'd like to also thank all of my colleagues in the scientific and medical spheres, whom I've found to be endlessly generous with their time and patience.

I'd like to thank my parents and my sister, Alice, for all your support over the past four years and the many years before. There's no way that any of this would have heppened without you all. Thank you.

Most of all, I want to express my deep gratitude to my newly minted wife, Kate. Sorry for all the late evenings and weekends spent working. Thank you so much for going through this with me. I love you.
\end{acknowledgements}

\setcounter{tocdepth}{2} 
% Setting this higher means you get contents entries for
%  more minor section headers.

\tableofcontents
\listoffigures
\listoftables

\cleardoublepage

\glsadd{sym:x}
\glsadd{sym:y}
\printglossary[type=\acronymtype,title={List of Abbreviations}]
\printglossary[type=symbols,title={Notation}]

\cleardoublepage
